%% abstract.tex
\begin{abstract}

The global power grid is undergoing a fundamental structural transformation:
synchronous machines—whose high-inertia fault signatures underpinned a century
of deterministic relay design—are being displaced at accelerating pace by
Inverter-Based Resources (IBRs). Grid-Following (GFL) and Grid-Forming (GFM)
converters are voltage-controlled current sources whose fault currents are
limited to 1.1--2.0~pu by fast sub-cycle current limiters, and whose positive-,
negative-, and zero-sequence contributions are determined by software control
laws rather than by physical machine reactance. This paradigm shift
systematically invalidates the foundational assumptions of legacy overcurrent
(ANSI~51/67), impedance distance (ANSI~21), and differential (ANSI~87)
protection algorithms.

This thesis proposes a hierarchical, self-adaptive protection and control
architecture—termed the \textbf{Self-Adaptive Multi-Bus Protection (SAMBP)}
framework—for IBR-dominated distribution and sub-transmission networks.
The framework spans six technical tiers, each developed from first principles
and validated through Monte Carlo simulation, digital twin emulation, and
hardware-in-the-loop testing.

\textbf{Line Differential Protection (Chapter~4)} develops a family of adaptive
87L algorithms. The physics-constrained adaptive threshold modulates the
operating slope $k_m$ as a function of the real-time IBR current ratio
$\kibr = I_\mathrm{IBR}/I_\mathrm{rated}$, eliminating the static
bias margin that previously caused missed trips under weak-infeed conditions.
Supplementary harmonic-discrimination (87LH), negative-sequence (87LN), and
transient-differential (TDCS-87L) elements extend sensitivity across the full
spectrum of unsymmetrical and high-impedance fault types.

\textbf{Distance Protection (Chapter~5)} derives the apparent impedance
trajectory seen by a distance relay when one or both line ends are IBR-fed.
The infeed-corrected reach equation and a Quadrilateral characteristic
adapted to IBR reactive current injection are proposed, and a Monte Carlo
framework quantifies the reliability gain over classical mho elements.

\textbf{IBR Microgrid Protection (Chapter~6)} characterises the protection-relevant
distinction between GFL and GFM inverter modes and proposes a ROCOF-supervised
mode-switch islanding detection scheme that achieves isolation within 120~ms
while remaining compliant with IEEE~1547-2018 ride-through requirements.
An adaptive overcurrent setting-group scheme uses IEC~61850~GOOSE mode
flags to reconfigure relay parameters within 50~ms of a grid-to-island transition.

\textbf{Digital Twin (Chapter~7)} introduces a non-blocking real-time digital twin
architecture that runs an Extended Kalman Filter (EKF) to jointly estimate
the IBR current ratio $\kibr$, line impedance $\Zline$, and fault location
parameter $\alpha$ from current measurements alone. The EKF achieves
steady-state RMSE of 0.024 in $\kibr$, 0.5\% in $\Zline$, and 1.1\% in $\alpha$
at a 50~$\mu$s update cycle.

\textbf{Physics-Informed Machine Learning (Chapter~8)} develops a Physics-Informed
Neural Network (PINN) fault classifier that embeds three physics constraints—current
ratio, sequence symmetry, and impedance consistency—directly in the loss function.
The PINN achieves 96.1\% five-class accuracy on synthetic IBR fault data while
guaranteeing zero physics-constraint violations. A CUSUM change-point detector
triggers selective online re-training in response to fleet composition drift,
recovering full accuracy within 21 decision cycles with no catastrophic forgetting.

\textbf{Centralised Protection and Control (Chapter~9)} integrates the foregoing
elements into a Wide-Area Protection Coordinator (WAPC) that localises faults
via PMU voltage-dip triangulation (94.5\% accuracy), optimises relay settings
through an IEC standard-inverse CTI-graded engine, secures GOOSE channels with
HMAC-SHA256 (0.36~ms overhead, 54.9\% risk reduction), and validates the complete
system through a 100-scenario hardware-in-the-loop regression suite.

Collectively, this thesis contributes 18 original algorithms and architectures,
of which 6 are Tier-1 (novelty~$= 5$, significance~$= 5$ on independent
expert assessment). The proposed SAMBP framework provides the first complete,
closed-loop, self-adaptive protection architecture for mixed GFL/GFM
power systems, from local relay physics through centralised coordination
and cybersecurity.

\textbf{Keywords:}
Inverter-Based Resources, Line Differential Protection, Distance Protection,
Microgrid Protection, Digital Twin, Physics-Informed Neural Network,
Fault Classification, Wide-Area Protection, IEC~61850 GOOSE Security,
Adaptive Relay Coordination.

\end{abstract}
